\documentclass[a4paper]{report}
\usepackage[utf8]{inputenc}
\usepackage[portuguese]{babel}
\usepackage{hyperref}
\usepackage{a4wide}
\hypersetup{pdftitle={LI4 - Proposta de Projeto},
pdfauthor={João Teixeira, José Ferreira, Miguel Solino, Maria Silva, Pedro
Oliveira},
colorlinks=true,
urlcolor=blue,
linkcolor=black}
\usepackage{subcaption}
\usepackage[cache=false]{minted}
\usepackage{listings}
\usepackage{booktabs}
\usepackage{multirow}
\usepackage{appendix}
\usepackage{tikz}
\usepackage{authblk}
\usepackage{bashful}
\usepackage{verbatim}
\usepackage{amsmath}
\usepackage{tikz}
\usepackage{tikz,fullpage}
\usepackage{pgfgantt}
\usetikzlibrary{arrows,%
                petri,%
                topaths}%
\usepackage{tkz-berge}
\usetikzlibrary{positioning,automata,decorations.markings}
\AfterEndEnvironment{figure}{\noindent\ignorespaces}
\AfterEndEnvironment{table}{\noindent\ignorespaces}

\begin{document}

\title{LI4 - Proposta de Projeto\\ 
\large PL1 - Grupo 1.1}
\author{José Ferreira (A83683) \and João Teixeira (A85504) \and Maria Silva
(A83840) \and Miguel Solino (A86435) \and Pedro Oliveira (A83762)}
\date{\today}

\begin{center}
    \begin{minipage}{0.75\linewidth}
        \centering
        \includegraphics[width=0.4\textwidth]{images/eng.jpeg}\par\vspace{1cm}
        \vspace{1.5cm}
        \href{https://www.uminho.pt/PT}
        {\color{black}{\scshape\LARGE Universidade do Minho}} \par
        \vspace{1cm}
        \href{https://www.di.uminho.pt/}
        {\color{black}{\scshape\Large Departamento de Informática}} \par
        \vspace{1.5cm}
        \maketitle
    \end{minipage}
\end{center}

\pagebreak
\section{Título e Enquadramento}
Um dos grandes problemas de gestão que todas as pessoas enfrentam ao longo da
sua vida é como gerir as suas compras de produtos alimentares. É necessário
gerir o que comprar, quando comprar e como comprar de forma a manter a dispensa
recheada. E os problemas não acabam por aí, depois de comprar os produtos é
necessário decidir o que fazer com eles e impedir que eles se estraguem. Com as
nossas vidas cada vez mais ocupadas e privadas de tempo, algo simples que junte
estas funcionalidades todas seria uma lufada de ar fresco.
The Kitchen App

\section{Motivações}
As motivações que nos levaram a iniciar este projeto foram o querer simplificar
o processo de ir as compras e de gestão das nossas dispensas de casa.
Queremos resolver problemas comuns do quotidiano, como comprar produtos que 
ja temos em demasia ou querer cozinhar aquele prato
que nos apetece mesmo e faltar algum dos ingredientes. 
Desta forma, ser-nos-á possível facilitar a vida aos utilizadores e, 
ao mesmo tempo, reduzir o desperdício alimentar.

\section{Objetivos}
O objetivo deste projeto é criar uma aplicação que facilite a gestão e
aproveitamento da despensa em casa permitindo também a criação de despensas
partilhadas que podem ser editadas por todos os residentes.
Para facilitar o aproveitamento dos ingredientes a aplicação irá sugerir várias
receitas que poderão posteriormente ser adicionadas aos favoritos. De forma a
diminuir o desperdício alimentar a aplicação também informará o utilizador
quando os produtos estão quase a ficar fora do prazo de validade.

\section{Requisitos}
\subsection{Requisitos Funcionais}
\begin{itemize}
    \item Criação/Gestão de Utilizador
    \item Criação/Gestão de despensas
    \item Adicionar utilizador a uma despensa
    \item Gerir produtos na despensa
    \item Gerir produtos
    \item Gerir listas de compras
    \item Procurar receitas (filtrar com base nos produtos disponíveis)
    \item Gerir receitas favoritas
\end{itemize}

\subsection{Requisitos não-funcionais}
\begin{itemize}
    \item Um utilizador deve ser criado fornecendo um nome, password, e-mail e
        data de nascimento
    \item Uma despensa deve ser criada com nome e convidados (opcional)
    \item Os produtos devem ser adicionados com o nome, tipo e data de validade
        (opcional)
    \item Na procura de produtos pode ser fornecido o tipo a pesquisar
    \item A criação de uma lista de compras pode ser feito com base numa receita
    \item As receitas pesquisadas podem ser ordenadas com base nos produtos
        disponíveis na dispensa
    \item Receitas podem ser adicionadas por utilizadores
    \item Um utilizador pode ter receitas favoritas
    \item Segurança
    \item Eficiência
\end{itemize}

\section{Tecnologias}
\begin{itemize}
    \item C\#
    \item MySQL
    \item JavaScript
    \item HTML/CSS
    \item Python
    \item Bash
\end{itemize}

\pagebreak
\section{Cronograma}
\noindent\resizebox{\textwidth}{!}{
\begin{tikzpicture}
\begin{ganttchart}[vgrid, hgrid]{1}{10}
    \ganttbar{Criação de Repositório}{1}{1} \\    
    \ganttbar{Definir Objetivos}{1}{2} \\    
    \ganttbar{Use Cases}{2}{2} \\    
    \ganttbar{Estudo de Mercado}{2}{2} \\
    \ganttbar{Modelação do Projeto}{3}{3} \\
    \ganttbar{Modelação de Base de Dados}{3}{4} \\
    \ganttbar{Planeamento da UI}{3}{4} \\
    \ganttbar{Desenvolvimento}{5}{10} \\
    \ganttbar{Escrita do Relatório}{1}{10} \\
    \gantttitle{1}{2}
    \gantttitle{2}{2}
    \gantttitle{3}{2}
    \gantttitle{4}{2}
    \gantttitle{5}{2}
\end{ganttchart}
\end{tikzpicture}
}

\end{document}
